% Source printed pp. 445-446 / PDF pp. 453-454.
\setlength{\emergencystretch}{4em}
\chapter{半稳定向量丛的模空间}\label{src:9}
\hypersetup{pageanchor=false}
\AddToHook{shipout/foreground}[chapter-nine-index-pages]{\put(0,0){\special{pdf:dest (page.\thepage) [@thispage /Fit]}}}
\begin{epigraph}
向量丛及其相应的模空间在代数几何中具有根本的重要性. 近几十年来, 这一课题与数学其他领域的联系进一步丰富了它的内涵. 这些领域包括微分几何 (它一直发挥着根本作用), 拓扑学 (模空间不仅具有有趣的拓扑特征, 而且有助于解决拓扑问题), 以及或许最令人惊讶的理论物理, 特别是规范理论, 量子场论和弦论.
\epigraphauthor{编者 \cite[前言]{BBPGPR09}}
\end{epigraph}
本章证明定理~\hyperref[thm:B]{B}: 对特征为 $0$ 的代数闭域 $\Bbbk$ 上亏格为 $g$ 的光滑连通射影曲线 $C$, 参数化秩为 $r$, 次数为 $d$ 的半稳定向量丛的叠 $\cBun_{r,d}^{\mathrm{ss}}(C)$ 是光滑, 普遍闭且不可约的代数叠, 维数为 $r^2(g-1)$, 并有射影良模空间 $M_{r,d}^{\mathrm{ps}}(C)$. 我们遵循 \S\bookref{1.7.2-subsection} 概述的构造射影模空间的六步策略:
\begin{enumerate}[(1)]
\item (代数性) 将 $\cBun_{r,d}^{\mathrm{ss}}(C)$ 表为 $\Coh(C)$ 的子叠; 后者参数化 $C$ 上所有凝聚层, 是局部有限型代数叠 (定理~\bookref{4.1.23}).
\item (半稳定性的开放性) $\cBun_{r,d}^{\mathrm{ss}}(C)\subseteq\Coh(C)$ 是开子叠 (定理~\bookref{9.2.4}).
\item (半稳定性的有界性) $\cBun_{r,d}^{\mathrm{ss}}(C)$ 是有限型的 (定理~\bookref{9.2.8}).
\item (半稳定约化) 由 Langton 定理 (\bookref{9.3.1}), $\cBun_{r,d}^{\mathrm{ss}}(C)$ 普遍闭.
\item (模空间的存在性) 因为 $\cBun_{r,d}^{\mathrm{ss}}(C)$ 是 $\Theta$-完备且 S-完备的 (命题~\bookref{9.3.8}), 良模空间存在定理 (\bookref{7.10.1}) 蕴含存在良模空间 $\cBun_{r,d}^{\mathrm{ss}}(C)\to M_{r,d}^{\mathrm{ps}}(C)$, 其中 $M_{r,d}^{\mathrm{ps}}(C)$ 是固有代数空间, 且 $\dim M_{r,d}^{\mathrm{ps}}(C)=r^2(g-1)+1$ (定理~\bookref{9.3.11}).
\item (射影性) 遵循 Faltings 的方法验证 $M_{r,d}^{\mathrm{ps}}(C)$ 是射影的.
\end{enumerate}
我们还在 \S\bookref{9.6} 给出 $M_{r,d}^{\mathrm{ps}}(C)$ 的 GIT 构造. 在秩 $r=1$ 的情形, $\cBun_{1,d}(C)$ 与参数化 $C$ 上次数为 $d$ 的线丛的 Picard 叠 $\underline{\mathcal{Pic}}_C^d$ 等同. 我们在定理~\bookref{7.4.65} 中已证明存在带 $\mathbb G_m$-带的 gerbe $\underline{\mathcal{Pic}}_C^d\to\underline{\mathrm{Pic}}^d(C)$ (特别地, 它是粗模空间), 其中 $\underline{\mathrm{Pic}}^d(C)$ 是射影概形. 若还有 $d=0$, 则 $\underline{\mathrm{Pic}}^0(C)$ 是 Abel 簇, 群律由张量积给出.

\textbf{记号:} 本章始终以 $C$ 表示域 $\Bbbk$ 上的光滑几何连通射影曲线. 我们常进一步假设 $\Bbbk$ 代数闭和/或特征为 $0$.\footnote{本章大多数结果在 $\Bbbk$ 为任意特征 $0$ 的域时也成立, 且证明相同. 特征 $0$ 假设仅在应用良模空间存在定理 (定理~\bookref{7.10.1}) 时是必要的.} 由于 $C$ 固定, 我们从模叠和模空间的记号中略去它, 例如 $\cBun_{r,d}^{\mathrm{ss}}:=\cBun_{r,d}^{\mathrm{ss}}(C)$ 和 $M_{r,d}^{\mathrm{ps}}:=M_{r,d}^{\mathrm{ps}}(C)$.
